\documentclass[11pt,a4paper]{article}
\usepackage[T1]{fontenc}
\usepackage{lmodern}
\usepackage[margin=27mm]{geometry}
\usepackage{amsmath,amssymb,amsthm,booktabs,microtype,xcolor,hyperref,fancyhdr}
\hypersetup{colorlinks=true,linkcolor=black,citecolor=blue,urlcolor=blue,pdfauthor={Jihun Kim},pdftitle={An Explicit 26-Measurement Phase-Retrieving Frame in C8: Version 2}}
\setlength{\emergencystretch}{3em}
\setlength{\parskip}{0.25em}
\setlength{\headheight}{14pt}
\pagestyle{fancy}\fancyhf{}
\fancyhead[L]{\small Explicit phase retrieval in $\mathbb C^8$}
\fancyhead[R]{\small Proof Candidate v2}\fancyfoot[C]{\thepage}
\newtheorem{theorem}{Theorem}
\newtheorem{lemma}[theorem]{Lemma}
\newtheorem{corollary}[theorem]{Corollary}
\newcommand{\CC}{\mathbb C}
\newcommand{\RR}{\mathbb R}
\newcommand{\norm}[1]{\left\lVert#1\right\rVert}
\newcommand{\co}[2]{[t^{#1}]\,#2}
\title{An Explicit 26-Measurement\\Phase-Retrieving Frame in $\mathbb C^8$\\[0.5em]
\large A five-parameter family and its exact boundary\\Proof Candidate -- Version 2}
\author{Jihun Kim\\\small Independent Researcher\\\small\texttt{jihunkimkw@gmail.com}}

\date{20 September 2026}
\begin{document}\maketitle
\begin{abstract}
We construct a five-parameter family of 26-measurement frames in $\CC^8$.
For real parameters $b,c,s,m,\gamma$ with $0<s<1$, $m<0$, and $\gamma>0$,
the family gives phase retrieval precisely when $c>b^2/s$.
The measurements determine squared moduli of polynomial evaluations at any
26 distinct positive real points. Three vanishing coefficients and a bound
on positive roots imply that these values determine the squared-modulus
polynomial. To prove
uniqueness up to phase, we compare the average and difference of two
possible signals. Coefficient elimination reduces the problem to positive
quadratic expressions and an invertible two-by-two matrix. When $c\le b^2/s$,
we construct nonproportional polynomials with identical squared moduli.
The family contains the previously published frame and a representative
with twenty monomial terms. The result gives $m_{\CC}(8)\le26$;
minimality is not established.
\end{abstract}
\noindent\textbf{Status.} Proof candidate; not peer reviewed.
Section~\ref{sec:status} describes
the completed checks and their scope.

\section{The question, and a short explicit answer}
A measurement of $x\in\CC^8$ is a real number $|q^*x|^2$, where
$q^*x=\sum_{k=0}^7\overline{q_k}x_k$.
Multiplying $x$ by a complex scalar of modulus one leaves every such
measurement unchanged. We ask whether 26 choices of $q$ distinguish
vectors in $\CC^8$ up to this unavoidable common phase.
A family with this property is called a phase-retrieving frame.
For any dimension $d$, let $m_{\CC}(d)$ denote the least number of such
measurements needed to distinguish all vectors in $\CC^d$ up to a common phase.

Consider the following eight polynomials:
\begin{equation}\label{eq:sparse}
\begin{aligned}
r_0&=1+it+it^5-it^{13},&
r_1&=t^2+\tfrac i2t^3-it^{11},\\
r_2&=t^4+it^5-it^9,&
r_3&=t^6-it^7,\\
r_4&=t^8+it^9+it^{13},&
r_5&=t^{10}+\tfrac i2t^{11},\\
r_6&=t^{12}+it^{13},&r_7&=t^{14}.
\end{aligned}
\end{equation}
For $j=1,\ldots,26$ take
\begin{equation}\label{eq:sparseframe}
 q_j=2\bigl(\overline{r_0(j)},\ldots,\overline{r_7(j)}\bigr)^T.
\end{equation}
Every entry is a Gaussian integer. The proof below shows that equality of
these 26 measurements forces two vectors to differ by a common phase.
The eight defining polynomials in \eqref{eq:sparse} contain twenty monomial
terms in total.

The geometric idea is to represent a signal by a polynomial $f$.
With our normalization, the measurement at $t_j$ is $4|f(t_j)|^2$
and therefore determines $|f(t_j)|^2$.
We arrange two properties: 26 samples determine this whole real polynomial,
and the polynomial determines $f$ up to phase. For the second property,
align their lowest-degree nonzero coefficients at $t=0$ and write them as
$v+h$ and $v-h$. Equal squared moduli mean that $v(t)$ and $h(t)$ are
perpendicular in the real plane $\CC$ at every real $t$.
The chosen space makes this possible only when $h=0$.

We work with a larger family to identify the sign condition that separates
uniqueness from ambiguity. The basis in \eqref{eq:sparse} is obtained from
the parameter choice $(b,c,s,m,\gamma)=(0,1,\tfrac12,-1,1)$ by an invertible
linear transformation, given in Section~\ref{sec:examples}. The frame
in the first release~\cite{KimV1} is another member of the family.

Conca, Edidin, Hering, and Vinzant proved that $4d-4$ generic complex
measurements suffice in dimension $d$~\cite{CEHV}. Vinzant constructed
eleven measurements in dimension four~\cite{Vinzant}. Here we give explicit
26-measurement frames in dimension eight and characterize the successful
parameters within our family.

\section{A family with a sharp boundary}\label{sec:family}
For a polynomial $f(t)=\sum c_jt^j$, put
$f^\#(t)=\sum\overline{c_j}t^j$. Thus $ff^\#$ has real coefficients and,
for real $t$, equals $|f(t)|^2$.
We write $[t^j]f$ for the coefficient of $t^j$ in $f$; for a matrix of
polynomials this notation is applied entry by entry. Conjugation also
acts entry by entry, and $a^{\#T}$ means $(a^\#)^T$.
Let $b,c\in\RR$ and assume
\begin{equation}\label{eq:domain}
 0<s<1,\qquad m<0,\qquad \gamma>0.
\end{equation}
Define
\begin{equation}\label{eq:seed}
 a=\begin{pmatrix}1+it+bt^3+ict^5\\t^2+ist^3-bt^5\\t^4+it^5\end{pmatrix},
 \qquad J=\begin{pmatrix}2c&0&2\\0&2s&0\\2&0&0\end{pmatrix}.
\end{equation}
Since $\det J=-8s$, the matrix is invertible on the whole domain, including
$c\le b^2/s$. Concatenate the following entries to form $p=(p_0,\ldots,p_7)^T$:
\begin{equation}\label{eq:assembly}
 p=\bigl(a+\gamma t^8J^{-1}a^\#,\quad t^6+imt^7,\quad t^8a,\quad t^{14}\bigr)^T.
\end{equation}
The transpose in this notation only indicates a column of eight entries.
Let $U=\operatorname{span}_{\CC}\{p_0,\ldots,p_7\}$.
Each $p_j$ begins with $t^{2j}$, with coefficient one, so these entries are
linearly independent and $\dim_{\CC}U=8$.

The identity $[t^5](aa^T)=iJ$ explains the choice of $J^{-1}$:
the conjugated correction cancels the degree-thirteen cross terms,
as shown in Section~\ref{sec:sampling}. Without this correction, the space
would contain both $a_0$ and $t^8a_0$, and hence the ambiguous pair
$a_0(1+it^8)$, $a_0(1-it^8)$.

\begin{theorem}\label{thm:main}
Under \eqref{eq:domain}, let $t_1,\ldots,t_{26}$ be any distinct positive real
numbers and set $q_j=2\overline{p(t_j)}$.
The vectors $q_j$ give phase retrieval if and only if
\begin{equation}\label{eq:sharp}
 \Delta:=c-\frac{b^2}{s}>0.
\end{equation}
When $\Delta\le0$, the space $U$ contains two nonproportional polynomials
with identical squared-modulus polynomials. Consequently, no set of real
sampling points can distinguish all elements of $U$ up to phase by their
squared moduli.
\end{theorem}

The quantity $\Delta$ is the minimum of the quadratic expression $R$ in
\eqref{eq:RD}. Positivity of $R$ is the key sign condition in the uniqueness
proof. Section~\ref{sec:failure} constructs ambiguous pairs when
$\Delta\le0$.

For $f_x=\sum x_kp_k$ one has $q_j^*x=2f_x(t_j)$.
Thus the measurement is $|q_j^*x|^2=x^*(q_jq_j^*)x$, with the rank-one
matrix $q_jq_j^*$. Every $q_j$ is nonzero since $p_7(t_j)=t_j^{14}>0$.
When the parameters and sampling points are rational, a common integer
scaling makes every entry of every measurement vector a Gaussian integer.

\section{Why 26 positive sampling points suffice}\label{sec:sampling}
All polynomials in $U$ have degree at most 14. For every $u\in U$,
\begin{equation}\label{eq:missing}
 [t]uu^\#=[t^5]uu^\#=[t^{13}]uu^\#=0.
\end{equation}
The seed satisfies
$[t](aa^{\#T})=[t^5](aa^{\#T})=0$ and $[t^5](aa^T)=iJ$;
these follow by multiplying the three displayed entries in \eqref{eq:seed}.
The correction starts in degree eight and hence does not affect the first
two missing degrees. In products $p_jp_k^\#$ with $j,k<3$, the correction
contributes the following matrix to the coefficient of degree thirteen:
\[
 \gamma\bigl(iJJ^{-T}-iJ^{-1}J\bigr)=0.
\]
Products between a corrected seed entry and an entry of $t^8a$, with
coefficient conjugation on one factor, reduce to the vanishing degree-five
coefficient of $aa^{\#T}$. The squared modulus of $p_3=t^6+imt^7$ is
$t^{12}+m^2t^{14}$. Products involving $p_3$ and an entry of
the corrected seed or $t^8a$ have no term of degree thirteen; all other
products start in higher degree. This proves \eqref{eq:missing}
for all products $p_jp_k^\#$ and hence for every $u$.

\begin{lemma}\label{lem:roots}
A nonzero real polynomial with at most $M$ nonzero monomials has at most
$M-1$ distinct positive roots.
\end{lemma}
\begin{proof}
Induct on the actual number $N\le M$ of nonzero monomials.
The case $N=1$ is immediate. For $N\ge2$, divide by the lowest power of
$t$, which preserves positive roots. The resulting polynomial has a
nonzero constant term and $N-1$ other monomials. Its derivative is nonzero
and has $N-1$ monomials, so by induction it has at most $N-2$ distinct
positive roots. Rolle's theorem gives at most $N-1\le M-1$ distinct
positive roots for the original polynomial.
\end{proof}

For $f,g\in U$, the real polynomial $ff^\#-gg^\#$ has degree at most 28.
Of the 29 possible coefficients, three vanish by \eqref{eq:missing}.
Thus it has at most 26 nonzero monomial terms. If 26 sampled squared moduli agree,
Lemma~\ref{lem:roots} gives
\begin{equation}\label{eq:normid}
 ff^\#=gg^\#.
\end{equation}
The use of distinct \emph{positive} nodes is essential to this argument.
It remains to understand when \eqref{eq:normid} forces a common phase.

\section{Comparing an average with its difference}\label{sec:normal}
If one polynomial in \eqref{eq:normid} is zero, so is the other.
For $f\ne0$, let $i$ be the least index with a nonzero coefficient in
its expansion in the basis $p_0,\ldots,p_7$. If the first term of $f$ at
zero is $\lambda t^{2i}$, the first term of $ff^\#$ is
$|\lambda|^2t^{4i}$. Equation~\eqref{eq:normid} therefore gives the same
initial index $i$ for $g$, with initial coefficient of modulus $|\lambda|$.
Rotate the phases of $f$ and $g$ independently and divide both by this common modulus so that
both initial coefficients become one. Then
\begin{equation}\label{eq:vh}
\begin{aligned}
 v&=(f+g)/2=p_i+\sum_{j>i}(x_j+iy_j)p_j,\\
 h&=(f-g)/2=\sum_{j>i}(r_j+i\alpha_j)p_j,\\
 F&=hv^\#+h^\#v=0.
\end{aligned}
\end{equation}
Since $ff^\#-gg^\#=2(hv^\#+h^\#v)$, equality of the squared-modulus
polynomials is equivalent to $F=0$. Here $x_j,y_j,r_j,\alpha_j\in\RR$
are coordinates of the normalized average and difference. For real $t$,
$F(t)=2\operatorname{Re}(h(t)\overline{v(t)})$, so $F=0$ expresses
pointwise orthogonality in the real plane $\CC$. The normalization divides
only by the nonzero initial coefficient modulus and imposes no restriction
on measurement zeros, common factors, or repeated roots.

For each fixed $v$, the map $h\mapsto hv^\#+h^\#v$ is real-linear on
$\operatorname{span}_{\CC}\{p_{i+1},\ldots,p_7\}$, viewed as a real vector
space. We prove its injectivity for every allowed $v$ by triangular
coefficient elimination, followed by scalar equations with strictly signed
factors and a two-dimensional linear system in the last case.

The expansions at $t=0$ are
\[
 p_j=t^{2j}(1+i\sigma_jt)+O(t^{2j+2}),\qquad
 (\sigma_0,\ldots,\sigma_7)=(1,s,1,m,1,s,1,0).
\]
If $h\ne0$ and $j>i$ is its first nonzero index, the first two possible
coefficients of $F$ are $2r_j$ and, after $r_j=0$,
$2(\sigma_i-\sigma_j)\alpha_j$. Higher indices cannot contribute to them.
It follows that $r_j=0$ and $\sigma_i=\sigma_j$.
The four values $1,s,m,0$ are distinct under \eqref{eq:domain}, so the only
possible pairs are
\begin{equation}\label{eq:pairs}
 (i;j)=(0;2,4,6),\ (1;5),\ (2;4,6),\ (4;6).
\end{equation}
For $i=3,5,6,7$ there is no possible nonzero difference.
Thus the argument also covers signals whose first several $p$-basis coordinates vanish.

Write $E_k=[t^k]F$. Eliminate the difference coefficients in increasing
order of $j$. The coefficient of $F$ at degree $2i+2j$ determines
$r_j$ by division by 2. After that substitution, the next coefficient
determines $\alpha_j$ by division by the fixed nonzero number
$2(\sigma_i-\sigma_j)$ whenever the slopes differ. Equal slopes leave
$\alpha_j$ free and the corresponding equation as a constraint.
The term involving $\alpha_j$ in its even-degree equation is zero because
the leading coefficients are real. No later index can affect these equations.
The recurrence uses only these fixed nonzero divisors.
Appendix~\ref{app:subs} records all substitutions needed below.

\section{The positive side of the boundary}\label{sec:suff}
Assume now $\Delta>0$. Put
\begin{equation}\label{eq:RD}
\begin{aligned}
 R(X,Y)&=c-2bY+s(1-s)X^2+sY^2\\
       &=\Delta+s(1-s)X^2+s(Y-b/s)^2>0,\\
 D&=c+\gamma-\frac{b^2}{s-m}
   =\Delta+\gamma+\frac{(-m)b^2}{s(s-m)}>0.
\end{aligned}
\end{equation}
These inequalities force the remaining difference coefficients in \eqref{eq:pairs} to vanish.
Every equation below is a coefficient of the same $F$ after the stated
preceding substitutions.

\subsection{The average starts with \texorpdfstring{$p_4$ or $p_1$}{p4 or p1}}
For $i=4$, the only possible first difference is at 6. Eliminating through
degree 23 gives
\[
 E_{25}=2\alpha_6R(x_5,y_5).
\]
Thus $\alpha_6=0$ and all remaining coefficients of $h$ vanish.
For $i=1$, start the difference at 5 and eliminate through degree 17. Then
\[
 E_{21}=-2\alpha_5\left[
 \gamma(1+x_2^2+y_2^2)+\frac{m(m-s)}s x_3^2-my_3^2\right].
\]
The bracket is at least $\gamma>0$, so again $h=0$.

\subsection{The average starts with \texorpdfstring{$p_2$}{p2}}
Start the difference at 4 and eliminate through degree 16. One obtains
\[
 E_{17}=-2\alpha_4\left[
 \frac{(1-m)(s-m)}{1-s}x_3^2
 +(s-m)\left(y_3-\frac b{s-m}\right)^2+D\right].
\]
The coefficients of the squared terms are positive, and $D>0$, so $\alpha_4=0$.
The eliminated coefficients before index 6 also become zero. Restarting
at 6 and eliminating through degree 19 yields
\[
 E_{21}=-2\alpha_6\left[\gamma+m(m-1)x_3^2-my_3^2\right].
\]
The bracket is at least $\gamma$, which forces $\alpha_6=0$ and $h=0$.

\subsection{The average starts with \texorpdfstring{$p_0$}{p0}: removing the first difference}
Start the difference at 2 and eliminate through degree 8. Then
\[
 E_9=2\alpha_2\left[
 \frac{(1-s)(s-m)}{1-m}x_1^2
 +(s-m)\left(y_1-\frac b{s-m}\right)^2+D\right].
\]
This forces $\alpha_2=0$, and every coefficient of $h$ before index 4
vanishes. The remaining case has two free imaginary coordinates.

\subsection{The final two-dimensional system}
For $i=0$ and $h$ starting at 4, degrees 8 through 15 give
\begin{equation}\label{eq:lastsubs}
\begin{aligned}
 r_4&=0,&r_5&=-\alpha_4y_1,&\alpha_5&=\alpha_4x_1,\\
 r_6&=-\alpha_4y_2,\\
 r_7&=\alpha_4x_2y_1-\alpha_4y_3-\alpha_6y_1,\\
 \alpha_7&=(1-m)\alpha_4x_3+(s-1)\alpha_4x_1x_2
                    +(1-s)\alpha_6x_1.
\end{aligned}
\end{equation}
Set
\begin{equation}\label{eq:RST}
\begin{aligned}
 \rho&=\alpha_6-\alpha_4x_2,\qquad R=R(x_1,y_1),\\
 T&=m\bigl((1-s)x_1x_3+y_1y_3\bigr)-\gamma x_2,\\
 S&=-\gamma(1+x_1^2+y_1^2+x_2^2+y_2^2)
       -m(m-1)x_3^2+my_3^2.
\end{aligned}
\end{equation}
Direct multiplication after \eqref{eq:lastsubs} gives
\begin{equation}\label{eq:matrix}
 \begin{gathered}
 E_{17}=2(R\rho+T\alpha_4),\qquad E_{21}=2(S\alpha_4+T\rho),\\
 \begin{pmatrix}S&T\\T&R\end{pmatrix}
 \begin{pmatrix}\alpha_4\\\rho\end{pmatrix}
 =\frac12\begin{pmatrix}E_{21}\\E_{17}\end{pmatrix}.
\end{gathered}
\end{equation}
The quadratic form is positive on the second coordinate axis and negative
on the first: $R\ge\Delta>0$, $S\le-\gamma<0$.
In two dimensions the off-diagonal term contributes $-T^2$ to the determinant,
so
\begin{equation}\label{eq:det}
 RS-T^2\le-\gamma\Delta<0.
\end{equation}
The matrix is therefore invertible for every average $v$.
Equivalently, with no division by its determinant,
\[
 R E_{21}-T E_{17}=2\alpha_4(RS-T^2).
\]
Since $E_{17}=E_{21}=0$, $\alpha_4=0$ and then $\rho=\alpha_6=0$.
Equation~\eqref{eq:lastsubs} makes every other difference coefficient zero.

All eight initial indices are now covered. Every solution of $F=0$
satisfies the coefficient equations used above, and these equations
force $h=0$.
Undoing the normalization shows that the original polynomials differ by a common unit phase. Together with
Section~\ref{sec:sampling}, this proves the sufficient direction of
Theorem~\ref{thm:main}.

\section{Ambiguous pairs on and beyond the boundary}\label{sec:failure}
The last four entries of \eqref{eq:assembly} span
$t^8U_4$, where
\[
 U_4=\operatorname{span}_{\CC}\{a_0,a_1,a_2,t^6\}.
\]
For $c\le b^2/s$, we seek $v,h\in U_4$ such that $hv^\#+h^\#v=0$
and $v+h$ is not proportional to $v-h$. Multiplying this pair by $t^8$
gives an ambiguous pair in $U$, independently of $m$ and $\gamma$.

Start at the boundary $c=b^2/s$. The polynomials
$v=a_0+i(b/s)a_1$ and $h=ia_2-(b/s)t^6$ belong to $U_4$ and simplify to
\[
 v=1+it+i\frac bs t^2,\qquad h=it^4v,
 \qquad f_\pm=v(1\pm it^4).
\]
On the real line, the factors $1+it^4$ and $1-it^4$ have identical
moduli. This gives an explicit ambiguity at the boundary.

To reach $c<b^2/s$, we adjust the average and difference while keeping
both in $U_4$. The coefficients $W,K,Z$ below are chosen so that every
coefficient of $hv^\#+h^\#v$, except those of degrees nine and twelve,
vanishes identically. We solve the remaining two equations on a real curve
in the $(X,Y)$-plane.
For real $X,Y$ put
\begin{align*}
 W&=b+sXY-sY,&K&=-bX+(s-1)WX-W,\\
 Z&=bsXY+c(1-s)X+WY,\\
 v&=a_0+(X+iY)a_1+iW a_2+(Z+iK)t^6,\\
 h&=ia_2+(-Y+i(1-s)X)t^6.
\end{align*}
Define
\[
\begin{aligned}
 R_0&=c-2bY+s(1-s)X^2+sY^2,\\
 G&=(b(2-s)+s^2Y)X+b-sY,\qquad H=(s-1)X-Y^2.
\end{aligned}
\]
Direct multiplication gives the full polynomial identity
\begin{equation}\label{eq:failureid}
 hv^\#+h^\#v=2R_0t^9+2\bigl(HG-(1-s)XYR_0\bigr)t^{12}.
\end{equation}
By \eqref{eq:failureid}, $R_0=G=0$ implies $hv^\#+h^\#v=0$.

If $b=0$, take $Y=0$ and $X=\sqrt{-c/(s(1-s))}$.
This is real for $c\le0$ and gives $R_0=G=0$.
If $b\ne0$, solving $G=0$ for $Y$ on $0\le X<1/s$ gives the curve
\[
 0\le X<1/s,\qquad
 Y=\frac{b((2-s)X+1)}{s(1-sX)}.
\]
Along this curve, $G=0$ and
\begin{equation}\label{eq:curveR}
 R_0(X)=c-\frac{b^2}s+s(1-s)X^2+
             \frac{4b^2X^2}{s(1-sX)^2}.
\end{equation}
For $c=b^2/s$ choose $X=0$. For $c<b^2/s$, this expression is negative
at zero, tends to $+\infty$ as $X\to(1/s)^-$, and has derivative
\[
 2s(1-s)X+\frac{8b^2X}{s(1-sX)^3}>0\qquad(0<X<1/s).
\]
Thus, when $b\ne0$ and $c<b^2/s$, there is a unique root in $0<X<1/s$.
If $b,c,s$ are rational, that root is specified by the quartic equation
\[
 \bigl(s^2(1-s)X^2+sc-b^2\bigr)(1-sX)^2+4b^2X^2=0
\]
together with the isolating interval $0<X<1/s$.

The resulting pair $v+h,v-h$ has the same squared-modulus polynomial by \eqref{eq:failureid}.
Both have constant coefficient one, whereas $h$ has first term $it^4$.
They cannot differ by a constant phase: the constant coefficients would
force that phase to be one, contradicting $h\ne0$.
Multiplying by $t^8$ proves the necessary direction of
Theorem~\ref{thm:main}.

\section{Two explicit frames}\label{sec:examples}
\subsection{Recovering the first released frame}
At
\[
 (b,c,s,m,\gamma)=(1,4,\tfrac12,-2,1)
\]
formula~\eqref{eq:assembly} gives the eight polynomials of the first
release~\cite{KimV1}. In particular $\Delta=2>0$.
The convention $q_j=2\overline{p(j)}$ is unchanged. Thus the family theorem
recovers the original frame as well as its phase-retrieval conclusion.
The exact archived input is retained with the verification data in
Section~\ref{sec:status}.

\subsection{The twenty-term representative}
At $(b,c,s,m,\gamma)=(0,1,\tfrac12,-1,1)$, one has $\Delta=1>0$.
From its assembly basis $p$ form
\[
 r_0=p_0-\tfrac12p_6,\quad r_1=p_1-p_5,\quad
 r_2=p_2-\tfrac12p_4+\tfrac12p_6,\quad r_j=p_j\ (3\le j\le7).
\]
This is a real triangular basis change with determinant one and gives exactly
\eqref{eq:sparse}. To check its effect on signals, write $r=Ap$.
Then $\sum x_kr_k=(A^Tx)^Tp$, and $A^T$ is invertible.
Uniqueness up to phase in the $p$ coordinates therefore implies the same
uniqueness in the $r$ coordinates. The resulting frame is different from the first
released one, even though both prove the same upper bound.

\begin{corollary}
The explicit Gaussian-integer frame \eqref{eq:sparseframe} gives
$m_{\CC}(8)\le26$.
\end{corollary}

For either successful example, an equivalent matrix statement is useful.
Let $\mathcal A_Q(H)=(q_j^*Hq_j)_{j=1}^{26}$ on Hermitian matrices.
The frame vectors $q_1,\ldots,q_{26}$ span $\CC^8$: a perpendicular vector would give a polynomial
of degree at most 14 vanishing at 26 distinct points, hence the zero polynomial.
If $H\ge0$ lies in the kernel, each $q_j^*Hq_j=0$ implies $Hq_j=0$,
so spanning gives $H=0$. The same holds for $H\le0$.
An indefinite rank-two Hermitian matrix is $xx^*-yy^*$ for nonproportional
$x,y$ by spectral decomposition, and the phase-retrieval property excludes it.
Thus
\[
 \ker\mathcal A_Q\cap\{H=H^*:0<\operatorname{rank}H\le2\}=\varnothing.
\]

\section{Scope and verification}\label{sec:status}
Theorem~\ref{thm:main} classifies the family \eqref{eq:assembly} in the
domain $0<s<1$, $m<0$, $\gamma>0$ and includes the original frame~\cite{KimV1}.
The measurement bound remains $m_{\CC}(8)\le26$. The minimum number of
measurements, including whether 25 suffice, is not determined here.
We make no assertion of generic phase retrieval with 26 measurements.
The positivity estimates used in the proof are expressed after
signal-dependent normalization and do not provide a global stability
constant or a noise-recovery guarantee.

The proof uses polynomial algebra, Rolle's theorem, and elementary linear
algebra.

Three scripts using SymPy check the coefficient identities, the
specializations, the basis transformation, and the frame entries.
The family and boundary scripts provide the first two check groups below.
The third script was transcribed separately and imports neither the original
family checker nor its basis data. Together the checks cover the three missing
degrees in all 64 ordered products $p_jp_k^\#$, all 28 pairs $0\le i<j\le7$,
the elimination and boundary identities, and all 416 entries across the
two frames. The execution logs record the following completed checks:
\begin{center}
\begin{tabular}{lr}
\toprule
Check group & Checks completed\\
\midrule
Original family checker &224\\
Boundary-identity checker &11\\
Separate transcription &847\\
\bottomrule
\end{tabular}
\end{center}
The inequalities and the logical arguments covering all cases, restarts,
and nonproportionality are given in the text. The exact frame data are in
\begin{itemize}
\item \texttt{data/v1\_frame\_gaussian\_integer.json} for the original example;
\item \texttt{data/v2\_sparse\_frame\_gaussian\_integer.json} for the twenty-term example.
\end{itemize}


The research, exposition, and computational checks were developed with
substantial AI assistance. External expert review and proof-assistant
verification have not been completed. The first release remains available at its DOI~\cite{KimV1}.

\appendix
\section{The coefficient substitutions}\label{app:subs}
The lists below record the substitutions used in Section~\ref{sec:suff}.
In each list, the $p$-basis coefficients of $h$ below the indicated index
vanish, while higher coefficients remain free. The recurrence in
Section~\ref{sec:normal} gives the displayed substitutions. All denominators
depend only on the fixed parameters and are nonzero under \eqref{eq:domain}.
The final case $(i,j)=(0,4)$ appears in \eqref{eq:lastsubs}.
\subsection*{$i=4$, first difference at 6: degrees 20 through 23}
\[
 r_6=0,\qquad r_7=-\alpha_6y_5,\qquad
 \alpha_7=(1-s)\alpha_6x_5.
\]
\subsection*{$i=1$, first difference at 5: degrees 12 through 17}
\[
\begin{aligned}
 r_5&=0,&r_6&=-\alpha_5y_2,&\alpha_6&=\alpha_5x_2,\\
 r_7&=-\alpha_5y_3,&\alpha_7&=\frac{s-m}{s}\alpha_5x_3.
\end{aligned}
\]
\subsection*{$i=2$, first difference at 4: degrees 12 through 16}
\[
\begin{aligned}
 r_4&=0,\qquad r_5=-\alpha_4y_3,\qquad
 \alpha_5=\frac{1-m}{1-s}\alpha_4x_3,\\
 r_6&=\alpha_4\left[-b+(s-m)y_3-\frac{s-m}{1-s}x_3y_3-y_4\right].
\end{aligned}
\]
\subsection*{$i=2$, restarted difference at 6: degrees 16 through 19}
\[
 r_6=0,\qquad r_7=-\alpha_6y_3,\qquad
 \alpha_7=(1-m)\alpha_6x_3.
\]
\subsection*{$i=0$, first difference at 2: degrees 4 through 8}
\[
\begin{aligned}
 r_2&=0,\qquad r_3=-\alpha_2y_1,\qquad
 \alpha_3=\frac{1-s}{1-m}\alpha_2x_1,\\
 r_4&=\alpha_2\left[b+(m-s)y_1+\frac{s-m}{1-m}x_1y_1-y_2\right].
\end{aligned}
\]
In the cases $(i,j)=(2,4)$ and $(0,2)$, these formulas leave the higher
coordinates free. If $\alpha_4=0$ in the first case, all difference
coefficients below index 6 vanish; if $\alpha_2=0$ in the second, all
coefficients below index 4 vanish. The argument can therefore restart at
those indices, as in Section~\ref{sec:suff}.


\begin{thebibliography}{9}
\bibitem{CEHV} A. Conca, D. Edidin, M. Hering, and C. Vinzant,
\emph{An algebraic characterization of injectivity in phase retrieval},
Applied and Computational Harmonic Analysis 38 (2015), 346--356.
\href{https://doi.org/10.1016/j.acha.2014.06.005}{doi:10.1016/j.acha.2014.06.005}.
\bibitem{Vinzant} C. Vinzant, \emph{A small frame and a certificate of its injectivity},
2015 International Conference on Sampling Theory and Applications, 197--200.
\href{https://arxiv.org/abs/1502.04656}{arXiv:1502.04656}.
\bibitem{KimV1} J. Kim, \emph{An Explicit 26-Measurement Phase-Retrieving Frame
in C8: Proof Candidate}, first release, 19 September 2026.
\href{https://doi.org/10.5281/zenodo.22847200}{doi:10.5281/zenodo.22847200}.
\end{thebibliography}
\end{document}
